\section{Open Questions}

Five concrete follow-up probes on the Friedrich et al.\ 2012 GLOBLE / GLDM
static formulation. Each is scoped to be free-endpoint feasible on a laptop or
mesh-available GPU; none require paid endpoints or new wet-lab data.

\subsection*{Q1. Does the $\beta$-vs-$\alpha$ anti-correlation hold at PIDE scale?}

\textbf{Basis.} Our 17-line subset gave Pearson $r(\alpha,\beta)=+0.655$,
opposite sign to the paper's headline claim. Herr 2014 shows $\varepsilon_i$
and $\varepsilon_c$ themselves strongly co-vary in this subset, masking the
predicted anti-correlation. The paper's original claim was anchored on the
full $\sim$150-line PIDE dataset.

\textbf{Next steps.} Pull the PIDE v3 database (Friedrich 2013 IJRB companion)
or re-digitise Friedrich 2012 Fig.~5. Refit
$(\varepsilon_i, \varepsilon_c)$ per line against clonogenic data using an
L-BFGS optimizer with sensible priors. Recompute Pearson and Spearman
correlations at full scale. Report whether the anti-correlation survives and
over which $\alpha$-range and dose-window it is strongest.

\subsection*{Q2. Sensitivity to Hi-C-informed loop-length distribution}

\textbf{Basis.} The paper fixes $N_L = 3000$ identical loops. Modern Hi-C
resolves loop domains at kbp-to-Mbp scale with a heavy-tailed length
distribution that varies between cell types. The clustered-DSB probability
$p_c$ is nonlinear in per-loop mean occupancy, so a length distribution
matters even when the total DNA and total DSB count are held fixed.

\textbf{Next steps.} Download a public Hi-C loop-domain call set
(e.g., HFFc6 or GM12878 from 4DN). Compute the empirical loop-length CDF.
Replace the single-$N_L$ Poisson kernel in \verb|globle_static.py| with a
size-weighted mixture over loop lengths. Re-run the RT112 dose-response and
compare against the uniform-loop baseline. One-at-a-time sensitivity on
$N_L$, mean-loop-length, and shape parameter.

\subsection*{Q3. Interaction with modern DSB-repair pathway modeling}

\textbf{Basis.} GLDM in its 2012 static form treats a clustered lesion as
irreversibly lethal with weight $\varepsilon_c$, folding all downstream
biology into two fitted per-cell constants. Repair-pathway choice
(NHEJ vs HR vs alt-EJ) is cell-cycle and lesion-complexity dependent; modern
repair-kinetics models (MEDRAS, DSBrepair, TRAX-CHEM) resolve this
explicitly.

\textbf{Next steps.} Couple the GLDM cluster-probability output $p_c(D)$
to a two-compartment ODE repair model with distinct rate constants for
isolated vs clustered breaks, and compare against the empirical
$(\varepsilon_i, \varepsilon_c)$ fits. Target: published RT112 or HSG
$\gamma$H2AX foci time-course. Test whether a mechanistic
(repair-rate $+$ misrepair-probability) parameterisation predicts
$(\varepsilon_i, \varepsilon_c)$ rather than absorbing them as free
constants.

\subsection*{Q4. Application to FLASH radiotherapy}

\textbf{Basis.} Static GLDM has no time axis. FLASH radiotherapy
($\gtrsim$40~Gy/s) shows dose-rate-dependent normal-tissue sparing,
hypothesized to involve transient oxygen depletion or radical recombination
on sub-second timescales. A model with no time axis cannot mechanistically
predict FLASH sparing, but the question is whether the $\varepsilon_c$
cluster-lethality weight is empirically dose-rate-invariant.

\textbf{Next steps.} Pull published FLASH vs conventional survival curves
for a matched cell line (zebrafish embryo, mouse intestinal crypt, or
clonogenic in vitro). Fit static GLDM under conventional dose rate to
extract $(\varepsilon_i, \varepsilon_c)$; then attempt to fit the FLASH
curve. If GLDM cannot fit both with a single parameter set, quantify the
residual and propose a minimal FLASH extension (e.g., dose-rate-dependent
effective $\varepsilon_c$). The kinetic-GLOBLE branch (Herr 2014) is the
natural bridge.

\subsection*{Q5. ML-driven cluster-severity scoring}

\textbf{Basis.} GLDM's binary classification collapses cluster complexity
into a single indicator (1 vs $\geq 2$ DSBs per loop). Track-structure MC
produces a full spectrum of cluster complexity that is known to correlate
with residual damage and misrepair probability. An ML scorer trained on
MC-derived cluster fingerprints could give a continuous severity weight.

\textbf{Next steps.} Use a public TOPAS-nBio or Geant4-DNA output (or
generate a small one on a free endpoint) to build a training set of
(cluster-fingerprint, empirical lethality-weight) pairs. Train a lightweight
gradient-boosted regressor or shallow MLP. Plug the resulting severity
function into the GLDM survival integral, replacing the binary
$\varepsilon_i / \varepsilon_c$. Compare against the 17-line clonogenic
panel used here, then extend to a proton or carbon-ion dataset to test the
radiation-quality generalisation the paper's downstream ion-GLOBLE work
targets.
